% !TeX root = ../初等数论笔记.tex
% 由主文档通过 \input 载入；请编译主文档。
\setcounter{chapter}{2} % 保留原书章号。
\chapter{素数与算术基本定理}
\label{ch:primes}

正整数可以反复拆成较小整数的乘积，直到不能继续拆分。这样的终点就是素数。本章要证明两件不同的事：分解总能进行到底，而且无论怎样分解，得到的素因子及其重数都相同。唯一分解使许多关于整数的命题转化为关于非负整数指数的命题，是后续整除计算的共同基础。

\section{素数、合数与素数无穷性}

\begin{definition}
大于 $1$ 且正因数只有 $1$ 和自身的整数称为\emph{素数（prime number）}；大于 $1$ 的其他整数称为合数（composite number）。整数 $1$ 不属于这两类。
\end{definition}

每个 $n>1$ 都有素因子：取 $n$ 的大于 $1$ 的最小正因数 $d$；若 $d$ 为合数，它有一个介于 $1$ 与 $d$ 之间的因数，这也是 $n$ 的因数，矛盾。

\begin{theorem}[素数无穷性]
素数有无穷多个。
\end{theorem}
\begin{proof}
若全部素数为 $p_1,\ldots,p_r$，令 $N=p_1\cdots p_r+1>1$。取 $N$ 的素因子 $q$，则 $q$ 应是某个 $p_i$，于是同时整除 $p_1\cdots p_r$ 和 $N$，从而整除 $1$，矛盾。
\end{proof}

这里并未断言 $N$ 本身为素数，只断言它具有新的素因子。

例如 $2\cdot3\cdot5\cdot7\cdot11\cdot13+1=30031=59\cdot509$。它虽是合数，两个因子却都不在原先的列表中，足以实现证明目的。理解一个构造时，应检查证明究竟需要“新数是素数”，还是只需要“新数有新的素因子”；前一个要求在这里既不成立，也无必要。

\section{欧几里得引理与唯一分解}

\begin{lemma}[欧几里得引理（Euclid's lemma）]
\label{lem:euclid-prime}
若 $p$ 为素数且 $p\mid ab$，则 $p\mid a$ 或 $p\mid b$。
\end{lemma}
\begin{proof}
若 $p\nmid a$，由素数定义有 $\gcd(p,a)=1$，应用推论\ref{cor:coprime-cancellation}得 $p\mid b$。
\end{proof}

\begin{theorem}[算术基本定理（fundamental theorem of arithmetic）]
\label{thm:unique-factorization}
每个整数 $n>1$ 都有唯一表示
\begin{equation}
 n=p_1^{\alpha_1}\cdots p_r^{\alpha_r},\qquad
 p_1<\cdots<p_r,\quad \alpha_i\in\mathbb Z_{>0},
 \label{eq:prime-factorization}
\end{equation}
其中 $p_1,\ldots,p_r$ 为素数。
\end{theorem}
\begin{proof}
对 $n$ 强归纳证明存在性。若 $n$ 为素数，已成立；否则 $n=ab$，其中 $1<a,b<n$，把 $a,b$ 的素数分解相乘即可。为证唯一性，将任意两种分解展开为素数的乘积。第一种的任一素因子，由欧几里得引理必整除第二种中的某个素数，因而等于它；消去这一对因子后继续，最终所有素因子及重数一致。
\end{proof}

对负整数，另加一个负号即可。整数 $1$ 对应空乘积，值约定为 $1$。把 $1$ 排除在素数之外，正是为了使分解中的因子个数不受任意插入 $1$ 的影响。

唯一性证明中最关键的是素数的消去性质，不能把它直接推广给任意整数。例如 $6\mid2\cdot3$，但 $6$ 既不整除 $2$，也不整除 $3$。此外，两种分解不可能在消去公共因子后出现“一边还有素数，另一边为 $1$”：若一边还有因子，其乘积至少为 $2$。这一点补足了逐次消去论证的终止情形。

\section{素因子指数与整除判定}

对素数 $p$ 和非零整数 $n$，记 $v_p(n)$ 为 $p$ 在 $|n|$ 的分解中的指数，称为 $p$ 进赋值（$p$-adic valuation）。不出现的素因子指数为 $0$。由唯一分解，非零 $a,b$ 满足
\[
 v_p(ab)=v_p(a)+v_p(b),\qquad
 a\mid b\ \Longleftrightarrow\ v_p(a)\le v_p(b)\ \text{对每个素数 }p.
\]
因此最大公因数与最小公倍数的指数分别取两者的最小值和最大值。

若 $a,b,a+b$ 均非零，写 $a=p^ru,b=p^sv$ 且 $p\nmid uv$。当 $r<s$ 时，$a+b=p^r(u+p^{s-r}v)$，括号内不被 $p$ 整除；故
\[
 v_p(a+b)\ge\min\{v_p(a),v_p(b)\},
\]
且两指数不同时等号成立。指数相同则可能严格增大，例如 $v_2(1+3)=2$。

\begin{example}[把整除要求转化为指数要求]
求最小正整数 $k$，使 $72k$ 为完全平方数；再求最小正整数 $t$，使 $72\mid t^2$。
\end{example}
\begin{solve}
$72=2^3\cdot3^2$。第一问要求 $3+v_2(k)$ 和 $2+v_3(k)$ 均为偶数，并且其他素因子的指数也为偶数。最小选择是 $v_2(k)=1$、其余指数为零，故 $k=2$，此时 $72k=144$。第二问要求 $2v_2(t)\ge3$、$2v_3(t)\ge2$，所以 $v_2(t)\ge2$、$v_3(t)\ge1$，最小 $t=12$。两个问题都得到平方数 $144$，但变量满足的是不同的指数条件。
\end{solve}

\begin{corollary}
固定 $k\ge2$。正整数 $n$ 是整数的 $k$ 次幂，当且仅当每个 $v_p(n)$ 都是 $k$ 的倍数。特别地，互素正整数的乘积为平方时，各因子均为平方。
\end{corollary}
\begin{proof}
若 $n=t^k$，则 $v_p(n)=kv_p(t)$；反之把各指数除以 $k$ 后构造 $t$。互素因子不共享素因子，故乘积中的偶指数分别属于各因子。
\end{proof}

\section{阶乘、二项式系数与筛法}

\begin{proposition}[阶乘中的素因子指数]
对素数 $p$ 与整数 $n\ge0$，
\begin{equation}
 v_p(n!)=\sum_{j\ge1}\left\lfloor\frac{n}{p^j}\right\rfloor.
 \label{eq:factorial-valuation}
\end{equation}
\end{proposition}
\begin{proof}
$n!$ 中每个 $p$ 的倍数至少贡献一次，每个 $p^2$ 的倍数再贡献一次，依此类推。若某数的指数为 $s$，它恰好被计数 $s$ 次。$p^j>n$ 后各项为零，故这里只涉及有限求和；$n=0$ 时两边均为零。
\end{proof}

对 $0\le k\le n$，二项式系数 $\binom nk=n!/[k!(n-k)!]$ 的素因子指数由三次应用式\eqref{eq:factorial-valuation}得到。例如 $100!$ 中 $5$ 的指数为 $20+4=24$，$2$ 的指数更大，故十进制末尾恰有 $24$ 个零。

\begin{example}[分层计数为何不会重复多算]
计算 $v_2(10!)$，并说明因子 $8$ 被计数三次的原因。
\end{example}
\begin{solve}
偶数 $2,4,6,8,10$ 各贡献一个 $2$；$4,8$ 各再贡献一个；$8$ 再贡献一个。因此总指数为 $5+2+1=8$。因子 $8=2^3$ 本来就应贡献三个 $2$，三次出现恰好记录它的三层因子。这个计数方式统计的是素因子的重数，统计对象并不是互不重复的整数。
\end{solve}

二项式系数的整性也可先从组合意义理解：从 $n$ 个不同对象中选出 $k$ 个，不考虑次序，共有 $\binom nk$ 种方法。于是阶乘商确实是整数，可以使用素因子指数判定其整除性。例如
\[
 v_2\!\left(\binom{10}{4}\right)
 =(5+2+1)-(2+1)-(3+1)=1,
\]
与 $\binom{10}{4}=210$ 相符。

若 $n$ 为合数，写 $n=ab$ 且 $1<a\le b$，则 $a\le\sqrt n$；$a$ 的一个素因子也不超过 $\sqrt n$。因此试除只需检查不超过平方根的素数。\emph{埃拉托色尼筛法（sieve of Eratosthenes）}用于列出 $N$ 以下的所有素数：从 $2$ 起，遇到尚未标记的 $p$ 就保留它，并标记 $p^2,p^2+p,\ldots\le N$；处理至 $p^2>N$ 后停止。较小的合数倍数已由更小素因子标记，所有合数最终都会被某个不超过其平方根的素因子筛去。

\section{习题}

唯一分解把整除关系变为各素数指数的比较：乘法对应相加，最大公因数取最小指数，最小公倍数取最大指数。使用阶乘公式时，计数对象是素因子的重数；讨论素性时，找不到小因子必须配合平方根上界才能得出结论。

\begin{enumerate}
 \item 分解 $360$；证明它的每个正约数唯一写成 $2^a3^b5^c$，其中 $0\le a\le3,0\le b\le2,0\le c\le1$。
 \item 计算 $v_2(100!)$ 与 $v_5\bigl(\binom{100}{50}\bigr)$。
 \item 若正整数 $a,b$ 互素且 $ab$ 为立方，证明 $a,b$ 都是立方。
 \item 证明对素数 $p$ 和 $1\le k<p$，$p\mid\binom pk$。
 \item\label{ex:prime-cube} 求最小正整数 $k$，使 $180k$ 为整数的立方。
 \item\label{ex:prime-odd-divisors} 证明正整数的正约数个数为奇数，当且仅当它为完全平方数。分别用约数配对与素因子指数证明。
 \item\label{ex:prime-rational-root} 证明若正整数 $n$ 的平方根为有理数，则 $n$ 为完全平方数。
 \item\label{ex:prime-binomial} 对任意素数 $p$，计算 $v_p\bigl(\binom{p^2}{p}\bigr)$，并说明结果对 $p=2$ 是否成立。
\end{enumerate}

